\documentclass[11pt,a4paper]{article}
\usepackage[margin=25mm]{geometry}
\usepackage{fontspec}
\setmainfont{Latin Modern Roman}
\setlength{\parindent}{0pt}
\setlength{\parskip}{0.55em}
\title{Mapa natal demonstrativo\\As doze casas}
\author{Segundo lote para revisão editorial}
\date{}
\newcommand{\leitura}[3]{\subsection{#1}\textit{#2}\par #3}
\newcommand{\subtituloleitura}[1]{\textbf{#1}\par\nopagebreak[4]}
\begin{document}
\maketitle
\section{Casas do mapa}

\leitura{Casa 1 · Presença e iniciativa}{Libra, 20°13\ensuremath{^{\prime}} · Regente: Vênus na casa 3 · Ocupante: Júpiter}{
\subtituloleitura{Você encontra seu lugar pela qualidade da conversa}
Ao entrar num ambiente, você percebe rapidamente o que aproxima as pessoas e quais assuntos pedem mais tato. Sua presença ganha força quando encontra uma pergunta que muda o nível da troca. Júpiter em Escorpião nessa casa dá profundidade a essa curiosidade: você pode começar com uma observação discreta e, em poucos minutos, chegar ao que realmente mobiliza alguém.

Sua maneira de participar também se constrói com o tempo. Vênus em Capricórnio na casa 3 favorece vínculos cultivados em conversas que continuam depois da primeira impressão. A tensão de Vênus com a Lua ajuda a entender por que uma interação que correu bem por fora ainda pode pedir tempo para você reconhecer o que sentiu. Ao retomar a conversa, dê lugar também à sua impressão.
}

\leitura{Casa 2 · Recursos e valores}{Escorpião, 22°09\ensuremath{^{\prime}} · Regente tradicional: Marte na casa 5 · Ocupante: Netuno}{
\subtituloleitura{Você reconhece valor quando consegue dar forma ao que imagina}
Sua relação com recursos passa pelo que está envolvido na troca: confiança, investimento pessoal e a liberdade de fazer algo que tenha sentido para você. Com Netuno em Escorpião nesta casa, o valor emocional de um trabalho às vezes pesa tanto quanto o preço. Isso aparece quando você aceita uma proposta porque vê potencial nela e só depois calcula a energia que vai exigir.

Marte em Peixes na casa 5 liga seus recursos à capacidade de criar e experimentar. O sextil de Marte com o Sol na casa 3 sugere um caminho concreto: transformar uma ideia criativa em algo que você consiga explicar, apresentar e negociar. Num projeto novo, definir logo o que será entregue e quanto tempo você pode dedicar ajuda a proteger o próprio entusiasmo.
}

\leitura{Casa 3 · Aprendizado e comunicação}{Sagitário, 19°37\ensuremath{^{\prime}} · Regente: Júpiter na casa 1 · Ocupantes: Sol e Vênus}{
\subtituloleitura{Sua curiosidade ganha força quando vira uma maneira própria de explicar}
Você tende a aprender por perguntas que abrem o assunto, até encontrar uma ideia que valha organizar e transmitir. Com Sol e Vênus em Capricórnio nesta casa, conversar bem envolve construir continuidade. Uma explicação clara, uma pergunta que retoma o fio da semana anterior ou uma anotação que outras pessoas conseguem usar são formas de tornar seu pensamento presente.

Júpiter em Escorpião na casa 1 liga essa curiosidade à impressão que você deixa. A quadratura de Júpiter com Mercúrio dá contexto ao contraste entre querer ir direto ao ponto e precisar conferir detalhes. Numa discussão, expor primeiro a pergunta e depois mostrar o raciocínio dá a quem ouve uma maneira de acompanhar a profundidade a que você quer chegar.
}

\leitura{Casa 4 · Vida privada e base familiar}{Capricórnio, 15°31\ensuremath{^{\prime}} · Regente: Saturno na casa 7 · Ocupantes: Mercúrio e Fortuna}{
\subtituloleitura{Sua base se fortalece quando os compromissos de casa ficam claros}
Você tende a se sentir mais à vontade em casa quando as responsabilidades têm forma. Mercúrio em Capricórnio e a Parte da Fortuna nesta casa dão relevo à capacidade de organizar a vida privada de um jeito que se sustente. O alívio aparece em coisas concretas: uma decisão familiar bem encaminhada ou uma tarefa que ganhou um lugar estável na rotina.

Saturno retrógrado em Touro na casa 7 liga essa estrutura aos compromissos com quem divide a vida com você. A oposição a Júpiter na casa 1 põe lado a lado a vontade de decidir por conta própria e o tempo necessário para construir um acordo. Dizer com clareza o que você consegue assumir torna o combinado mais confiável para todos.
}

\leitura{Casa 5 · Criação, prazer e expressão}{Aquário, 13°01\ensuremath{^{\prime}} · Regente tradicional: Saturno na casa 7 · Ocupante: Marte}{
\subtituloleitura{Você cria melhor quando pode experimentar dentro de um compromisso claro}
Seu prazer em criar aparece quando existe liberdade para testar uma ideia incomum e ver aonde ela leva. Marte em Peixes nesta casa acrescenta um impulso que responde ao clima do momento: uma música, uma imagem ou a sugestão de alguém pode fazer um projeto avançar de repente. Em colaboração, você ganha espaço para improvisar quando os termos do trabalho já estão claros.

Saturno em Touro na casa 7 dá peso aos compromissos que tornam essa liberdade sustentável. O trígono de Saturno com Vênus na casa 3 favorece, na leitura, uma prática simples: conversar cedo sobre o que cada pessoa vai fazer. Um horário combinado ou uma divisão clara de tarefas deixa a energia criativa disponível para o projeto.
}

\leitura{Casa 6 · Rotina e trabalho cotidiano}{Peixes, 14°41\ensuremath{^{\prime}} · Regente tradicional: Júpiter na casa 1 · Ocupantes: Nodo Norte, Quíron e Vertex}{
\subtituloleitura{Sua rotina funciona melhor quando comporta iniciativa e ajuste}
Você trabalha com mais consistência quando consegue ajustar o plano às necessidades do dia. O Nodo Norte em Peixes põe essa flexibilidade em evidência. Quíron e Vertex em Áries dão outro tom à mesma rotina: há momentos em que a situação só anda quando você faz a primeira proposta, mesmo com parte do caminho em aberto.

Júpiter em Escorpião na casa 1 liga o funcionamento da rotina à sua disposição de investigar o que causa um problema. A oposição a Saturno na casa 7 acrescenta a importância dos acordos com quem trabalha ao seu lado. Se uma tarefa emperra toda semana, vale descobrir onde ela para e combinar um teste pequeno com as pessoas envolvidas. A solução ganha força quando cabe no dia real.
}

\leitura{Casa 7 · Relações e parcerias}{Áries, 20°13\ensuremath{^{\prime}} · Regente: Marte na casa 5 · Ocupante: Saturno}{
\subtituloleitura{Você precisa de franqueza e constância para confiar numa parceria}
Uma relação próxima ganha vida quando as preferências aparecem claramente e há espaço para fazer coisas juntos. Marte em Peixes na casa 5 leva essa iniciativa para o terreno da imaginação, da intimidade e do prazer de inventar algo a dois. O sextil com o Sol na casa 3 liga esse movimento à conversa: muitas vezes, a aproximação começa quando alguém põe em palavras uma vontade que já estava no ar.

Saturno retrógrado em Touro nesta casa dá peso ao que pode ser sustentado depois do entusiasmo inicial. Você tende a levar a sério aquilo que promete e a reparar se a outra pessoa consegue fazer o mesmo. Quando um acordo precisa mudar, ser direto desde cedo ajuda a preservar a confiança que os dois construíram.
}

\leitura{Casa 8 · Intimidade e recursos compartilhados}{Touro, 22°09\ensuremath{^{\prime}} · Regente: Vênus na casa 3}{
\subtituloleitura{Você aprofunda a confiança quando consegue falar sobre o que é compartilhado}
Na intimidade e nas decisões que envolvem recursos de mais de uma pessoa, você se orienta pelo que tem continuidade. Um compromisso se torna mais confiável quando existe clareza sobre o que cada um oferece, espera e pode sustentar. Vênus em Capricórnio na casa 3 leva esse tema para a conversa cotidiana: a confiança cresce conforme assuntos difíceis encontram palavras e os combinados são retomados quando necessário.

O trígono de Vênus com Saturno na casa 7 aproxima essa disposição da responsabilidade nas parcerias. Quando a decisão envolve uma despesa em comum ou uma responsabilidade delicada, definir como vocês vão revisar o acordo dá estabilidade à troca. Assim, mudanças de circunstância encontram um lugar para ser discutidas.
}

\leitura{Casa 9 · Visão de mundo e estudos}{Gêmeos, 19°37\ensuremath{^{\prime}} · Regente: Mercúrio na casa 4}{
\subtituloleitura{Você amplia sua visão ao pôr suas referências em conversa}
Você tende a se interessar por uma ideia nova quando consegue compará-la com algo que já conhece. Um estudo, uma viagem ou o contato com outro modo de viver pode tornar visível um pressuposto que veio da família e que, até então, parecia evidente. Com Mercúrio em Capricórnio na casa 4, essas descobertas ganham peso quando você consegue relacioná-las a uma experiência concreta.

A quadratura de Mercúrio com Júpiter na casa 1 dá contexto à vontade de tirar uma conclusão ampla a partir de uma descoberta importante. Você pode aproveitar esse impulso e ainda dar tempo para que outros relatos entrem na conversa. Uma boa pergunta, nesse caso, é o que muda na sua leitura quando alguém começa de uma história diferente da sua.
}

\leitura{Casa 10 · Trabalho público e direção}{Câncer, 15°31\ensuremath{^{\prime}} · Regente: Lua na casa 12 · Ocupante: Lilith}{
\subtituloleitura{Sua contribuição pública ganha forma quando você dá nome ao cuidado que oferece}
Você tende a perceber cedo o que uma equipe precisa para funcionar: alguém que receba quem está chegando, uma conversa que ficou pendente ou uma tarefa cujo peso ninguém calculou. A Lua em Libra na casa 12 liga essa percepção ao que você elabora longe da cena. A quadratura entre Lua e MC dá contexto ao esforço de conciliar a exposição pública com o tempo de que você precisa para recuperar seu equilíbrio.

Lilith em Leão nesta casa traz a autoria para o centro da questão. Quando o trabalho ganha visibilidade, vale conseguir dizer que decisão foi sua e qual efeito ela teve. A contribuição que você faz nos bastidores também merece aparecer na forma como responsabilidades e reconhecimento são distribuídos.
}

\leitura{Casa 11 · Amizades e projetos coletivos}{Leão, 13°01\ensuremath{^{\prime}} · Regente: Sol na casa 3}{
\subtituloleitura{Você ganha espaço num grupo quando ajuda uma ideia a avançar}
Em amizades e projetos coletivos, você tende a se envolver quando existe algo a construir e as pessoas têm abertura para contribuir com suas próprias ideias. O Sol em Capricórnio na casa 3 liga essa participação à maneira como você organiza uma conversa. Resumir uma decisão ou colocar uma proposta em termos que todos consigam usar é uma forma de liderança que nasce do trabalho compartilhado.

A conjunção do Sol com Vênus favorece a leitura desse talento pelo lado dos vínculos que ele cria. Você consegue reunir pessoas em torno de um assunto e cultivar a continuidade da troca. Se já fez a parte de estruturar a ideia, convidar outra pessoa a assumir o próximo passo ajuda o projeto a ganhar vida para além da sua iniciativa.
}

\leitura{Casa 12 · Recolhimento e elaboração}{Virgem, 14°41\ensuremath{^{\prime}} · Regente: Mercúrio na casa 4 · Ocupantes: Lua, Urano e Plutão}{
\subtituloleitura{É no recolhimento que você percebe quais detalhes ainda pedem atenção}
Depois de um dia cheio, você pode precisar de silêncio para entender por que uma conversa ou decisão continuou ocupando sua cabeça. A Lua e Urano em Libra trazem para esse espaço a revisão do que aconteceu nas relações. Plutão retrógrado em Virgem aprofunda o exame dos hábitos envolvidos, como a tendência a ajustar sempre a mesma parte de uma tarefa ou de uma convivência.

Mercúrio em Capricórnio na casa 4 liga essa elaboração à segurança de ter uma base privada onde você possa organizar o que percebeu. O trígono de Mercúrio com Plutão favorece a leitura de que dar nome a um padrão já ajuda a encontrar o detalhe que vale mudar. Depois de entender o que se repete, escolha uma mudança pequena que possa observar na prática.
}

\end{document}